\begin{theorem}
Let $Q$ be a type and let $r$ be a quasi-order relation on $Q$.  With
$\mathsf{HerCtblPower}(Q)$ the hereditary countable non-empty part of the
full tagged hierarchy, equipped with the restricted relation
$\mathsf{HerCtblRel}(r)$, one has
\[
\mathsf{Bqo}(r)\quad\Longleftrightarrow\quad
\mathsf{WellQuasiOrdered}(\mathsf{HerCtblRel}(r)).
\]
\end{theorem}
\begin{proof}
For the forward implication, assume $\mathsf{Bqo}(r)$.  By
\noderef{PowerQPresBqo}, the full tagged hierarchy with relation
$\mathsf{PowerRel}(r)$ is better-quasi-ordered.  Applying
\noderef{BqoImpliesWqo} to this relation gives
\[
\mathsf{WellQuasiOrdered}(\mathsf{PowerRel}(r)).
\]
Let $a:\mathbb N\to\mathsf{HerCtblPower}(Q)$ be arbitrary.  Regard it as
a sequence in the full hierarchy by setting
$b(k)=(a(k))._1$.  The displayed well-quasi-order statement supplies
$m<n$ with
\[
\mathsf{PowerRel}(r)(b(m),b(n)).
\]
By the definition of $b$ and by the restriction definition
\noderef{HerCtblRel}, this is exactly
\[
\mathsf{HerCtblRel}(r)(a(m),a(n)).
\]
Since $a$ was arbitrary, the hereditary countable carrier is
well-quasi-ordered.

Conversely, assume
$\mathsf{WellQuasiOrdered}(\mathsf{HerCtblRel}(r))$.  We prove
$\mathsf{Bqo}(r)$ using its definition in \noderef{Bqo}.  Let
$h:\mathsf{IncSeq}\to Q$ be locally constant and suppose that $h$ is
bad for $r$.  The nontrivial-front reduction
\noderef{BadMultiToNontrivialSuper}, applied to $r$, $h$, its local
constancy, and its badness, supplies a super-sequence
whose front is $\mathsf{DecidingFrontSet}(h)$ and whose values lie in $Q$.
This super-sequence is bad and
$\emptyset\notin\mathsf{DecidingFrontSet}(h)$.

For each $k\in\mathbb N$, define an element of the hereditary carrier by
\[
a(k)=\left\langle
  \widetilde f(\{k\}),
  \mathsf{HereditarilyCountable}(\widetilde f(\{k\}))
\right\rangle.
\]
This is well-defined as follows.  The front is nontrivial, so
\noderef{FrontTreeSingleton} gives
$\{k\}\in\mathsf{PrefixTree}(\mathsf{DecidingFrontSet}(h))$.
Consequently the singleton definition \noderef{TildeFSingleton} gives
the underlying value $\widetilde f(\{k\})$, and
\noderef{TildeFHerCtbl} supplies its hereditary-countability witness.
By the subtype definition \noderef{HerCtblPower}, the displayed pair is an
element of $\mathsf{HerCtblPower}(Q)$.

Apply the assumed hereditary well-quasi-ordering to this sequence.  It
gives $m<n$ such that
\[
\mathsf{HerCtblRel}(r)(a(m),a(n)).
\]
But \noderef{ConverseGame}, applied to the same bad super-sequence and
the fact that its front does not contain the empty set, says that for
every $m<n$ the singleton sequence has
\[
\neg\mathsf{HerCtblRel}(r)(a(m),a(n)).
\]
This contradicts the relation obtained from well-quasi-ordering.
Therefore no locally constant bad multi-sequence exists, and the
definition \noderef{Bqo} yields $\mathsf{Bqo}(r)$.
\end{proof}
